% Part: model-theory
% Chapter: interpolation
% Section: introduction

\documentclass[../../../include/open-logic-section]{subfiles}

\begin{document}

\olfileid[fr]{mod}{int}{int}
\olsection{Introduction}

Le théorème d'interpolation affirme ceci :
pour des !!{sentence}s $!A$ et $!B$, si $\Entails !A\lif !B$, il existe
!!a{sentence}~$!C$ tel que
$\Entails !A\lif !C$ et
$\Entails !C\lif !B$.
De plus, chaque !!{constant},
!!{function} et !!{predicate}
autre que~$\eq$ qui figure dans~$!C$
figure à la fois dans~$!A$ et dans~$!B$.
On appelle~$!C$ un \emph{interpolant}
de~$!A$ et~$!B$.

Le théorème présente un intérêt en lui-même.
Il permet aussi d'établir des résultats
sur la définissabilité dans une théorie
et sur les conditions auxquelles la réunion
de deux théories cohérentes reste cohérente.
Le premier de ces résultats est le
théorème de définissabilité de Beth;
le second, le théorème de satisfaisabilité
disjointe de Robinson.

\end{document}
